\subsection{ソフトウェアセキュリティツール}
セキュリティツールは、弱点を見つける、攻撃経路を確認する、開発者に早くフィードバックするために使われる。ただし、ツールは万能ではない。発見できる弱点はツールの方法と設定に依存し、結果の確認には専門知識が必要である。

\subsubsection{セキュリティ脆弱性確認ツール}
セキュリティ脆弱性確認ツールには、ソースコード解析器やバイナリ解析ツールがある。ソースコード解析器は、注入の欠陥、バッファオーバーフロー、安全でないライブラリ使用など、コードの書き方や論理に関係する弱点を探す。

バイナリ解析ツールは、コンパイル済みコードや第三者ライブラリを調べ、ソースコードでは見えにくい脆弱性や、コンパイル過程で生じる問題を探す。これらのツールは有用だが、発生条件が非常に特殊な状態や例外的な状況で現れる脆弱性まですべて見つけることはできない。

\subsubsection{ペネトレーションテストツール}
ペネトレーションテストツールは、運用環境またはそれに近い環境で、制御された攻撃を行い、脆弱性や弱点を見つけるために使われる。代表的な技法には、異常な入力やランダムな入力を与えるファジングがある。

ペネトレーションテストの価値は、実際の攻撃者がどのようにシステムを悪用し得るかについて洞察を得られる点にある。ただし、ツールを動かすだけでは不十分である。範囲、許可、ログ取得、影響の管理、結果の解釈が必要である。

\begin{notebox}{注意}ペネトレーションテストは必ず明示的な許可を得て行う。許可のない診断、侵入試行、スキャンは、学習目的であっても相手に被害を与え、法的問題を引き起こす可能性がある。\end{notebox}
